\section{Empirical strategy}
\label{sec:strategy}

\subsection{Local projections}

I estimate the response of house prices to the monetary policy shock with
panel local projections in the spirit of \citet{jorda2005}. The idea is to
run a separate regression for every horizon $h$, with the cumulative house
price change from $t-1$ to $t+h$ as the outcome. Compared to a VAR, this
approach does not force the impulse response into the shape implied by a
small dynamic system, and it makes it simple to interact the shock with a
country characteristic, which is exactly what my question requires.

The unit of observation is a country $c$ in quarter $t$. For each horizon
$h = 0, 1, \ldots, 12$ I estimate

\begin{equation}
\label{eq:lp}
\begin{split}
100 \left( \ln HP_{c,t+h} - \ln HP_{c,t-1} \right)
= {} & \beta_h \, \varepsilon_t
+ \delta_h \left( \varepsilon_t \times ARM_c \right)
+ \sum_{j=1}^{2} \phi_{h,j} \, \varepsilon_{t-j}
+ \sum_{j=1}^{2} \rho_{h,j} \, g_{c,t-j} \\
& + \gamma_h' X_{c,t-1}
+ \kappa_h \, C_t
+ \alpha_c
+ u_{c,t+h},
\end{split}
\end{equation}

where $HP_{c,t}$ is the real house price index, $\varepsilon_t$ is the
information cleaned monetary policy shock in basis points, $ARM_c$ is the
indicator equal to one if country $c$ has a variable rate dominant mortgage
market, $g_{c,t}$ is quarter on quarter real house price growth in percent,
$X_{c,t-1}$ contains lagged GDP growth and lagged inflation, $C_t$ is a dummy
for the pandemic quarters (2020Q1 to 2021Q2), $\alpha_c$ is a country fixed
effect, and $u_{c,t+h}$ collects everything else that moves house prices at
horizon $h$. Because the outcome is a log change times one hundred, the
coefficients read as the percent change in real house prices after a one
basis point tightening surprise.

Equation \eqref{eq:lp} is an estimating equation, not a structural model.
The parameter of interest is $\delta_h$, the difference in the house price
response between variable rate and fixed rate countries at horizon $h$. The
budget channel described in Section \ref{sec:mechanism} predicts
$\delta_h < 0$. In countries where mortgage payments react quickly to policy
rates, a tightening should depress house prices by more. Because the
interaction is included, $\beta_h$ is the response of the fixed rate
countries, the baseline group, and not the euro area average. It and the
remaining coefficients are auxiliary. The main effect of $ARM_c$ does not
appear in the equation because it is constant over time and therefore
absorbed by the country fixed effect. The baseline specification contains
no time fixed effects, because the shock is common to all countries and
time effects would absorb it completely. This is a real cost of the
design, and the robustness check in Section \ref{sec:timefe} addresses it
by dropping the main effect and keeping only the interaction.

\subsection{Identification}

Identification of the shock comes from the high frequency event window. The
surprise is measured as the change in the three month OIS rate within a
narrow window around the policy announcement. The identifying assumption is
that within this window, the only relevant news is the announcement itself,
so the surprise is unrelated to economic developments that were already
public before the meeting. This assumption is standard in the literature and
I take it from \citet{altavilla2019}. The sign restriction of
\citet{jarocinski2020} additionally removes events where the surprise mainly
revealed the central bank's private information about the economy. All
controls enter with a lag, so they are predetermined with respect to the
shock in quarter $t$.

For $\delta_h$ a second assumption is needed. The mortgage market
classification must not itself respond to monetary policy over the sample.
This is plausible, because the fixed versus variable character of a national
mortgage system is a slow moving institutional feature that depends on
funding models and regulation. The more serious concern is that $ARM_c$ is a single
binary trait of a country, so $\delta_h$ picks up the effect of anything
that is correlated with mortgage structure across countries, for example
supply constraints, homeownership rates, or the general boom and bust
history of the periphery. With eleven countries I cannot separate these
stories, and I interpret $\delta_h$ as a descriptive difference between the
two country groups rather than the pure causal effect of mortgage structure.

\subsection{Specifications and inference}

The baseline includes the standard local projection dynamics, two lags of
the shock and two lags of house price growth, because this is common
practice and it protects the estimate against serial correlation in shocks
and outcomes. This choice was made before comparing results across
specifications. I report the same equation without the dynamic terms as an
alternative, since the interaction turns out to be sensitive to this choice,
and I consider this sensitivity an informative result in itself. The lags
enter without an $ARM_c$ interaction, because interacting the full lag set
would add several parameters to a specification that is already demanding
for a sample of eleven countries. I therefore keep the interaction to the
contemporaneous term.

Standard errors are Driscoll Kraay \citep{driscoll1998} throughout. The
shock is the same for all countries in a given quarter, so the errors are
cross sectionally dependent by construction, and Driscoll Kraay standard
errors are built for exactly this situation. The estimation sample at $h = 0$ contains 1{,}151
country quarter observations for eleven countries from 1999Q4 to 2025Q4, and
it shrinks mechanically with the horizon because the outcome reaches $h$
quarters into the future.
